\chapter{沿着主线回看：每次增加了什么机制？}
从产生--降解出发，单节点正反馈通过多个稳态提供记忆，双节点互抑制通过放大表达差异形成开关，三节点负反馈增加动态环节而能使带摆动的小偏差从衰减转为增长，快慢前馈则以支路竞争形成瞬态响应。转向分子计数后，同一个长期平均值还对应一整个涨落分布；回到真实网络时，组合逻辑与模块耦合进一步决定这些局部机制能否被观察到。下表把各步得到的结论与它们的适用边界放在一起。
\begin{center}\begin{tabular}{>{\raggedright\arraybackslash}p{3.1cm}>{\raggedright\arraybackslash}p{6cm}>{\raggedright\arraybackslash}p{5.6cm}}\toprule
模型 & 由推导得到的行为 & 迁移结论时需注意\\\midrule
线性产生--降解 & 唯一指数稳定的稳态；出生--死亡离散模型为泊松稳态 & 泊松结论依赖恒定出生与一阶死亡\\
一维单调自抑制 & 唯一稳定的稳态；不能自行维持振荡 & 输入固定且没有延迟时成立\\
一维自激活 & 足够非线性时可能双稳态 & 多交点还需足够非线性与合适参数\\
双组分互相调控 & 正/负回路符号影响特征值和稳定性 & 稳定性仍依赖增益及具体产生率\\
三节点负反馈 & 小偏差可从衰减摆动转为增长摆动 & 稳定周期轨道需要非线性分析\\
前馈回路 & 可实现脉冲、滤波和特定条件下适应 & 精确适应还需稳态支路增益抵消\\
\bottomrule\end{tabular}\end{center}

分析一个新回路时，也可以沿这条路线：先从反应规则建立产生率，说明哪些快变量被省去；再寻找稳态，检查偏差会恢复还是放大；最后考虑分子涨落与其余网络的影响，判断模型能解释哪些实验现象。
